\section{Worum es geht}

\textbf{[SETZUNG]} Dieses Dokument behandelt den Casimir-Effekt im Rahmen der FFGFT.
Der Casimir-Effekt ist im Altkorpus die einzige experimentell zugängliche Brücke
zwischen der fundamentalen Längenskala $L_0=\xi\cdot\ell_P$ und messbaren Vakuumkräften.
Das Dokument zeigt: (1) Die modifizierte Casimir-Formel reproduziert die Standardformel
exakt -- keine neue Physik bei bisher zugänglichen Abständen. (2) Der Vorfaktor $4/3$
in $\xi=4/3\times10^{-4}$ ist durch den Casimir-Effekt unabhängig von Massenfits
bestätigt (A015). (3) Es gibt eine charakteristische Vakuum-Längenskala
$L_\xi\approx100\,\mu\text{m}$, bei der Casimir-Energiedichte und eine bestimmte
Vakuumgröße zusammenfallen -- das ist hier explizit
ausgeklammert.

\section{Fundamentale Längenskalen}

\textbf{[B]} FFGFT definiert eine Hierarchie charakteristischer Längenskalen:
\[
  L_0 = \xi\cdot\ell_P \approx 2{,}155\times10^{-39}\,\text{m}
  \quad(\text{Sub-Planck, minimale Raumzeitgranulation}),
\]
\[
  \ell_P = \sqrt{\hbar G/c^3} \approx 1{,}616\times10^{-35}\,\text{m}
  \quad(\text{Planck-Länge}),
\]
\[
  L_\xi \approx 100\,\mu\text{m}
  \quad(\text{charakteristische Vakuum-Längenskala, messbar}).
\]
Die Kopplungskonstante $\xi=4/3\times10^{-4}$ ist derselbe Parameter wie in A020.
Sie kann aus einem Lagrangian abgeleitet werden, der die Dynamik des Zeitfelds
beschreibt (A142).

\textbf{[B]} \textbf{Granulation der Raumzeit.} Bei $L_0$ sind alle Vakuumfluktuationen
vollständig wirksam. Für $d>L_0$ werden nur Teile dieser Kräfte sichtbar -- das ist
der physikalische Ursprung der $1/d^4$-Abhängigkeit der Casimir-Kraft. Aufgrund der
extrem kleinen Größe von $L_0\approx2\times10^{-39}\,\text{m}$ ist eine direkte
Messung nicht möglich.

\section{Die modifizierte Casimir-Formel und ihre Konsistenz}

\textbf{[K]} Die modifizierte Casimir-Energiedichte im FFGFT-Rahmen:
\[
  |\rho_\text{Casimir}(d)|
  = \frac{\pi^2}{240\,\xi}\,\rho_\text{vac}\cdot\Bigl(\frac{L_\xi}{d}\Bigr)^4,
\]
wobei $\rho_\text{vac}=\xi\hbar c/L_\xi^4$ eine charakteristische Vakuumdichte ist.
Einsetzen ergibt:
\[
  |\rho_\text{Casimir}(d)|
  = \frac{\pi^2}{240\,\xi}\cdot\frac{\xi\hbar c}{L_\xi^4}\cdot\frac{L_\xi^4}{d^4}
  = \frac{\pi^2\hbar c}{240\,d^4}.
\]
Das ist exakt die Standard-Casimir-Formel. Kein Widerspruch, kein neuer Parameter --
die FFGFT-Formulierung ist zu den etablierten QED-Ergebnissen numerisch äquivalent.

\textbf{[K]} \textbf{Numerische Verifikation} (Casimir-Energiedichte vs. Plattenabstand):

\begin{center}
\renewcommand{\arraystretch}{1.3}
{\small
\begin{tabular}{lcc}
\toprule
Abstand $d$ & $\rho_\text{Casimir}$ (J/m³) & Verhältnis zu $\rho_\text{vac}$ \\
\midrule
$100\,\mu$m & $4{,}17\times10^{-14}$ & 1{,}00 \\
$10\,\mu$m  & $4{,}17\times10^{-10}$ & $10^4$ \\
$1\,\mu$m   & $4{,}17\times10^{-2}$  & $10^{12}$ \\
\bottomrule
\end{tabular}
}
\renewcommand{\arraystretch}{1.0}
\end{center}

\textbf{[K]} \textbf{Charakteristische Skalen-Hierarchie:}
\[
  \frac{L_0}{\ell_P} = \xi = 1{,}333\times10^{-4},
  \qquad
  \frac{\ell_P}{L_\xi} \approx 1{,}6\times10^{-31}.
\]

\section{Bestätigung des Faktors 4/3}

\textbf{[B]} Der Vorfaktor $4/3$ in $\xi=4/3\times10^{-4}$ tritt unabhängig im
Casimir-Kontext auf: der Faktor $4/3$ erscheint im Kasimir-Energie-Term als
$E_\text{Casimir}\propto\tfrac43$ (Quarte -- musikalisches Intervall $4{:}3$) und ist
damit nicht nur eine Masse-Fit-Größe. Diese Unabhängigkeit ist in A015 als Casimir-Bestätigung
geführt: kein anderer Vorfaktor (Quinte $3/2$, Terz $5/4$) liefert konsistente
Ergebnisse für das Massenspektrum \emph{und} die Casimir-Struktur gleichzeitig.

\section{Experimentelle Tests}

\textbf{[K]} Drei Klassen von Tests:

\textbf{Präzisionsmessungen bei $d\approx L_\xi$:} Bei Plattenabstand $\approx100\,\mu$m
erreicht die Casimir-Energiedichte die Größenordnung von $\rho_\text{vac}$. Dieser
Bereich ist experimentell zugänglich.

\textbf{Skalierungsverhalten:} Die $1/d^4$-Abhängigkeit sollte bis in den
Nanometerbereich exakt erfüllt sein. Abweichungen bei $d\lesssim10\,\text{nm}$ könnten
auf die Raumzeitgranulation hindeuten.

\textbf{Experimentelle $L_\xi$-Werte aus der Literatur:} Parallele Platten:
$\sim228\,\text{nm}$; Kugel-Platte-Geometrie: $\sim1{,}75\,\mu\text{m}$ bis
$\sim18\,\mu\text{m}$. Die Streuung spiegelt geometrische Unterschiede wider
($F\propto1/d^4$ vs.\ $F\propto1/d^3$).

\section{Prüfskript}

\begin{itemize}[leftmargin=2em,itemsep=2pt,topsep=2pt]
  \item Numerische Verifikation der modifizierten Casimir-Formel (Reproduktion
    der Standardformel), Längenskalen-Hierarchie, Vergleich mit $L_\xi$-Werten:
    \texttt{python/A\_Serie\_Skripte/a260\_casimir.py}
\end{itemize}

\section{Was offen bleibt}

\textbf{Direkte Messung von $L_0$ ist nicht möglich.} Die minimale Längenskala
$L_0\approx2\times10^{-39}\,\text{m}$ liegt weit jenseits jeder experimentellen
Reichweite. Die charakteristische Skala $L_\xi$ ist experimentell zugänglich,
aber ihre genaue Verbindung zu $L_0$ setzt die FFGFT-Rahmenstruktur voraus.

\textbf{Abweichungen vom $1/d^4$-Gesetz bei $d\sim10\,\text{nm}$ sind nicht
gemessen.} Die Prognose einer Skalenänderung nahe der Sub-Planck-Grenze ist
qualitativ, quantitativ nicht ausgeführt.

\section{Ersetzt}

Dieses Dokument nimmt auf: Dok.~091 (modifizierte Casimir-Theorie, Längenskalen-Hierarchie,
numerische Verifikation, experimentelle Vorhersagen und Tests). \emph{Nicht}
übernommen: die Casimir-Vakuum-Skalierung aus Dok.~091 (ausgeschlossener Sektor,
außerhalb A-Serien-Kern) und die makroskopischen Vakuum-Implikationen.

\section{Verwendung von KI-Werkzeugen}

Dieses Dokument ist unter Verwendung KI-gestützter Werkzeuge entstanden. Verantwortung
für Inhalt und Fehler liegt beim Autor. Wortlaut der Regeln in A010.
